\documentclass[10pt,a4paper]{amsart}
\usepackage[margin=19mm]{geometry}
\usepackage{amsmath,amssymb,mathtools}
\usepackage[scr=rsfs,cal=euler]{mathalfa}
\usepackage{xfrac}
\usepackage[unicode,hidelinks,bookmarks=false]{hyperref}
\newcommand{\ie}{i.e.\ }
\newcommand{\cf}{cf.\ }
\newcommand{\resp}{respectively\ }
\newcommand{\Subsection}{\subsection}
\providecommand{\nobreakdash}{\nobreak\hbox{-}}
\setlength{\emergencystretch}{2em}
\multlinegap0pt
\usepackage{bm}
\usepackage{bbm}
\usepackage{dsfont}
\usepackage{xspace}
\usepackage{cancel}
\usepackage{mathtools}
\usepackage{amsthm}
\makeatletter
\@ifundefined{thm}{\newtheorem{thm}{Thm}}{}
\@ifundefined{lem}{\newtheorem{lem}[thm]{Lemma}}{}
\@ifundefined{prop}{\newtheorem{prop}[thm]{Proposition}}{}
\makeatother
\DeclareMathOperator{\sotimes}{\square}
\newcommand{\CAlg}{\mathrm{CAlg}}
\newcommand{\Fil}{\mathrm{Fil}}
\newcommand{\HR}{\mathrm{HR}}
\newcommand{\Sp}{\mathrm{Sp}}
\newcommand{\THR}{\mathrm{THR}}
\newcommand{\Z}{\mathbb{Z}}
\newcommand{\rD}{\mathrm{D}}
\newcommand{\rH}{\mathrm{H}}
\newcommand{\ul}{\underline}
\renewcommand{\L}{\mathbb{L}}
\title[Real Topological Hochschild Homology of Perfectoid Rings]{Real Topological Hochschild Homology of Perfectoid Rings}
\author{Jens Hornbostel, Doosung Park}
\date{Source: arXiv:2310.11183v3 (CC BY 4.0)}
\begin{document}
\maketitle

\noindent\emph{Source and licence.} Real Topological Hochschild Homology of Perfectoid Rings. Version arXiv:2310.11183v3 (CC BY 4.0), distributed under the Creative Commons Attribution 4.0 International licence, as recorded in the arXiv licence metadata for this exact version and shown on the abstract page https://arxiv.org/abs/2310.11183v3. Full rights notice: licenses/arxiv-2310.11183-CC-BY-4.0.txt.\par\smallskip

\noindent\textit{Editorial scope.} Counted unit: the Lem whose source label is \texttt{lem:perfectoid basechange} in the article (source0.tex lines 3675--3687, source label \texttt{lem:perfectoid basechange}), delivered together with the article's own complete original proof of it (source0.tex lines 3688--3749, one \texttt{proof} environment of 62 lines). No mathematics is written, repaired or re-derived: statement and proof are the article's own text line for line. Required context retained uncounted: limit and smash (lines 2316--2332); prop:hrcompleted (lines 2422--2434); lem:pseudo-coherence (lines 3498--3502); perfectoid basechange HR (lines 3584--3596). Every definition, display and label of those lines is retained unchanged. Omitted from this package and not redistributed: 1--2315, 2333--2421, 2435--3497, 3503--3583, 3597--3674, 3750--4435. Nothing inside a retained excerpt is omitted or altered: every retained body is delivered exactly as the article's lines are written, so no omission receipt of any kind accompanies this package. Citations: the retained excerpts carry no citation command at all, so no bibliography entry of the article is transcribed and no reference is redistributed. Reference adaptation: none was needed and none was made; every reference inside the delivered unit resolves inside it, so no unresolved reference or citation is produced. Printed slips kept literal: no printed word, symbol, hypothesis or number of the counted statement or of its proof was altered. Layout and class: the article's own class is \texttt{amsart} with packages including amsmath, amsthm, amssymb, tikz-cd, xcolor, hyperref, todonotes; the wrapper is the lane's pinned \texttt{amsart} editorial wrapper at 10pt on A4, which restates the theorem-like environments the retained text needs and adds the article's own macro definitions that the retained text needs (\texttt{\textbackslash CAlg}, \texttt{\textbackslash Fil}, \texttt{\textbackslash HR}, \texttt{\textbackslash L}, \texttt{\textbackslash Sp}, \texttt{\textbackslash THR}, \texttt{\textbackslash Z}, \texttt{\textbackslash rD}, \texttt{\textbackslash rH}, \texttt{\textbackslash sotimes}, \texttt{\textbackslash ul}), with extra wrapper packages (bm, bbm, dsfont, xspace, cancel, mathtools). The delivered numbering is the unit's own. Assets: no retained span contains a figure environment, a graphics-inclusion command, a PDF-page inclusion, a table or a TikZ picture; no image file is copied into this package, the package declares no asset dependency and no image file is redistributed with it. Licence: version arXiv:2310.11183v3 of this article is distributed under the Creative Commons Attribution 4.0 International licence, as recorded in the arXiv licence metadata for this exact version and shown on the abstract page https://arxiv.org/abs/2310.11183v3. A full rights notice is delivered at licenses/arxiv-2310.11183-CC-BY-4.0.txt. Characters: every retained excerpt is pure ASCII, so the delivered engine, XeLaTeX with its default Latin Modern text font, reports no missing character.
\begin{lem}
\label{limit and smash}
Let $A \in \CAlg^{\Z/2}$ be a commutative ring object of $\Sp^{\Z/2}$,
let $L$ and $M$ an $A$-modules,
and let $\cdots \to \Fil_1 M\to \Fil_0 M:=M$ be a complete filtration on $M$.
Assume that $A$, $L$, and $M$ are $(-1)$-connected and $\Fil_n M$ is $(n-1)$-connected for every integer $n$.
Then there are natural equivalence of $A$-modules
\[
\lim_n (L\wedge_A \Fil_n M)
\simeq
0
\text{ and }
\lim_n (L\wedge_A \Fil^n M)
\simeq
L\wedge_A \lim_n \Fil^n M.
\]
\end{lem}

\begin{prop}
\label{prop:hrcompleted}
Let $A$ be a commutative ring.
Then the natural morphism
\[
\THR(A;\Z_p)
\sotimes_{\THR(\Z)}^\L
\ul{\Z}
\to
\HR(A;\Z_p)
\]
in $\rD(\ul{A})$ is an equivalence.
\end{prop}

\begin{lem}
\label{lem:pseudo-coherence}
Let $R$ be a perfectoid ring.
Then $\HR(R;\Z_p)$ and $\THR(R;\Z_p)$ are pseudo-coherent complexes of $\ul{R}$-modules.
\end{lem}

\begin{lem}
\label{perfectoid basechange HR}
Let $R\to R'$ be a homomorphism of perfectoid rings.
Then the natural morphism
\[
\HR(R;\Z_p)
\sotimes_{\ul{R}}^\L
\ul{R'}
\to
\HR(R';\Z_p)
\]
in $\rD(\ul{R'})$ is an equivalence.
\end{lem}

\begin{lem}
\label{lem:perfectoid basechange}
Let $R\to R'$ be a homomorphism of perfectoid rings.
Then the natural morphism
\[
\THR(R;\Z_p)
\sotimes_{\ul{R}}^\L
\ul{R'}
\to
\THR(R';\Z_p)
\]
in $\rD(\ul{R'})$ is an equivalence.
\end{lem}

\begin{proof}
It suffices to show 
that the induced map of $\Z/2$-spectra
\[
\THR(R;\Z_p)
\wedge_{\rH \ul{R}}
\rH \ul{R'}
\to
\THR(R';\Z_p)
\]
is an equivalence.
By Proposition \ref{prop:hrcompleted} and Lemma \ref{perfectoid basechange HR},
the induced map of $\Z/2$-spectra
\[
\rH \ul{\Z}
\wedge_{\THR(\Z)}
\THR(R;\Z_p)
\wedge_{\rH \ul{R}}
\rH \ul{R'}
\to
\rH \ul{\Z}
\wedge_{\THR(\Z)}
\THR(R';\Z_p)
\]
is an equivalence.
As in the proof of Lemma \ref{lem:pseudo-coherence},
this implies that the induced map of $\Z/2$-spectra
\begin{align*}
&
\ul{\pi}_n(\THR(\Z))
\wedge_{\THR(\Z)}
\THR(R;\Z_p)
\wedge_{\rH \ul{R}}
\rH \ul{R'}
\\
\to
&
\ul{\pi}_n(\THR(\Z))
\wedge_{\THR(\Z)}
\THR(R';\Z_p)
\end{align*}
is an equivalence for every integer $n$.
By induction on $n$,
we see that the induced map of $\Z/2$-spectra
\begin{align*}
&
\tau_{\leq n}(\THR(\Z))
\wedge_{\THR(\Z)}
\THR(R;\Z_p)
\wedge_{\rH \ul{R}}
\rH \ul{R'}
\\
\to
&
\tau_{\leq n}(\THR(\Z))
\wedge_{\THR(\Z)}
\THR(R';\Z_p)
\end{align*}
is an equivalence for every integer $n$.
Take $\lim_n$ on both sides,
and use Lemma \ref{limit and smash} to conclude.
\end{proof}

\end{document}
